\documentclass[../../../include/open-logic-section]{subfiles}

\begin{document}

\olfileid[id]{sth}{spine}{valpha}

\olsection{Mendefinisikan Tahap-Tahap sebagai $V_\alpha$}

Dalam \olref[sth][ordinals][]{chap}, kita mendefinisikan pengurutan baik dan
ordinal (von Neumann). Dalam bab ini, kita akan menggunakannya untuk
mencirikan hierarki himpunan \emph{itu sendiri}. Untuk melakukannya, ingat
bahwa dalam \olref[sth][ordinals][opps]{sec}, kita mendefinisikan gagasan
ordinal suksesor dan limit. Kita menggunakan gagasan-gagasan tersebut dalam
definisi berikut:

\begin{defn}\ollabel{defValphas}
	\begin{align*}
	V_\emptyset &\defis \emptyset\\
	V_{\ordsucc{\alpha}} &\defis \Pow{V_\alpha} & &
	\text{untuk sembarang ordinal }\alpha\\
	V_{\alpha} &\defis \bigcup_{\gamma < \alpha} V_\gamma & &
	\text{ketika }\alpha\text{ merupakan ordinal limit}
\end{align*}
\end{defn}
\noindent
Ini akan menjadi definisi melalui \emph{rekursi transfinit} pada ordinal.
Dalam hal ini, kita perlu membandingkannya dengan definisi rekursif
fungsi-fungsi pada bilangan asli.\footnote{Bdk.\ definisi penjumlahan,
perkalian, dan eksponensiasi dalam
\olref[sfr][infinite][dedekind]{sec}.} Seperti ketika menangani bilangan
asli, kita mendefinisikan suatu kasus dasar dan kasus-kasus suksesor; tetapi
ketika menangani ordinal, kita juga perlu menjelaskan perilaku kasus-kasus
\emph{limit}.

Definisi $V_\alpha$ ini akan menjadi suatu tonggak penting. Secara informal,
kita telah memotivasi hierarki himpunan sebagai pembentukan himpunan melalui
\emph{tahap-tahap}. Pada dasarnya, $V_\alpha$ tepat merupakan tahap-tahap
tersebut. Namun, yang penting, ini merupakan pencirian \emph{internal} atas
tahap-tahap. Alih-alih menyarankan suatu \emph{model} yang mungkin bagi
teori tersebut, kita akan mendefinisikan tahap-tahap itu \emph{di dalam}
teori himpunan kita.

\end{document}
